% =============================================================
% Glossar-Einträge.
%
% Eingebunden in der Präambel (vor \begin{document}). Im Text
% verwendet via \gls{schluessel} an der ersten Erwähnung; weitere
% Auftritte können auch \gls{} sein, der Glossar-Eintrag wird dann
% nur noch normal gesetzt.
%
% Schlüssel sind kurz und sprechend; Anzeigename steht in `name=`.
% =============================================================

\newglossaryentry{modulo}{
  name=Modulo,
  description={Operation \(a \bmod n\): der Rest, der bei der Division
    von \(a\) durch \(n\) bleibt. Universalwerkzeug der
    Codierungstheorie}}

\newglossaryentry{paritaet}{
  name=Parität,
  description={Eigenschaft einer Bitfolge, die sich aus der Anzahl der
    Einsen ergibt: \emph{gerade Parität}, wenn die Anzahl gerade
    ist, sonst \emph{ungerade}. Eine Paritätsprüfung erkennt jede
    ungerade Anzahl gekippter Bits}}

\newglossaryentry{paritaetsbit}{
  name=Paritätsbit,
  description={Zusätzliches Bit, das so gewählt wird, dass die
    Gesamt-Parität (meist gerade) zustande kommt. Erkennt
    Einzel-Bitfehler}}

\newglossaryentry{redundanz}{
  name=Redundanz,
  description={Daten-Anteil, der über die reine Information
    hinausgeht und der Erkennung oder Korrektur von
    Übertragungsfehlern dient. Preis für Robustheit}}

\newglossaryentry{ean13}{
  name=EAN-13,
  description={European Article Number, 13-stelliger Strichcode auf
    fast jedem Konsumgut. Letzte Ziffer ist eine Prüfziffer (Modulo
    10 mit Gewichten 1 und 3)}}

\newglossaryentry{isbn10}{
  name=ISBN-10,
  description={Bis 2007 verwendete 10-stellige Buchnummer. Prüfziffer
    Modulo 11; deshalb gelegentlich „X" als Prüfsymbol nötig}}

\newglossaryentry{luhn}{
  name=Luhn-Algorithmus,
  description={Prüfziffer-Verfahren für Kreditkartennummern (Hans Peter
    Luhn, 1954). Verdoppelt jede zweite Ziffer von rechts und summiert
    die Quersummen}}

\newglossaryentry{hammingdistanz}{
  name=Hamming-Distanz,
  description={Anzahl der Stellen, an denen sich zwei gleich lange
    Wörter unterscheiden. Zentrales Maß für Codes}}

\newglossaryentry{mindestdistanz}{
  name=Mindestdistanz,
  description={Kleinste Hamming-Distanz zwischen je zwei Codewörtern
    eines Codes. Bestimmt die Korrektur-/Erkennungs-Fähigkeit:
    \(d-1\) Fehler erkennbar, \(\lfloor (d-1)/2 \rfloor\)
    korrigierbar}}

\newglossaryentry{code}{
  name=Code,
  description={Menge gültiger Codewörter. Die Hamming-Distanzen
    zwischen ihnen bestimmen die Korrektureigenschaften}}

\newglossaryentry{codewort}{
  name=Codewort,
  description={Element eines Codes, also eine gültig codierte Bit-
    oder Symbolfolge}}

\newglossaryentry{hammingschranke}{
  name=Hamming-Schranke,
  description={Untere Grenze für die Anzahl Bits, die zur
    Einzelfehler-Korrektur nötig sind: \(2^k \cdot (n+1) \le 2^n\).
    Wird vom \((7,4)\)-Hamming-Code mit Gleichheit erreicht („perfekter
    Code")}}

\newglossaryentry{hammingcode}{
  name=Hamming-Code,
  description={Familie linearer Codes (Richard Hamming, 1950), die
    einen Einzel-Bitfehler korrigieren. Der \((7,4)\)-Variante codiert
    4 Datenbits in 7 Codebits}}

\newglossaryentry{secded}{
  name=SECDED,
  description={Single Error Correction, Double Error Detection.
    Erweiterung eines Hamming-Codes um ein Gesamt-Paritätsbit:
    erkennt Doppelfehler zusätzlich zur Einzelfehler-Korrektur. In
    ECC-RAM Standard}}

\newglossaryentry{buendelfehler}{
  name=Bündelfehler,
  description={Mehrere benachbarte Bitfehler, etwa durch Kratzer auf
    einer CD oder Funkfading. Realer als zufällige Einzelfehler;
    motiviert Interleaving}}

\newglossaryentry{interleaving}{
  name=Interleaving,
  description={Verschachtelung von Codewörtern vor der Übertragung,
    sodass ein Bündelfehler aus jedem Codewort nur ein einziges Bit
    trifft. Steckt in jeder CD und vielen Funkstandards}}

\newglossaryentry{crc}{
  name=CRC,
  description={Cyclic Redundancy Check. Standardverfahren der
    Fehlererkennung. Daten werden als Polynom interpretiert und
    modulo eines Generatorpolynoms gerechnet}}

\newglossaryentry{generatorpolynom}{
  name=Generatorpolynom,
  description={Festes Polynom \(G(x)\) bei CRC oder Reed-Solomon, das
    die Code-Eigenschaften bestimmt. Codewörter sind durch \(G(x)\)
    teilbar}}

\newglossaryentry{koerper}{
  name=Körper,
  description={Mathematische Struktur, in der man uneingeschränkt
    addieren, subtrahieren, multiplizieren und dividieren kann
    (Ausnahme: Division durch Null). Beispiele: \(\mathbb{R}\),
    \(\mathbb{Q}\), \(\mathbb{Z}/p\mathbb{Z}\) für \(p\) prim}}

\newglossaryentry{endlicherkoerper}{
  name=Endlicher Körper,
  description={Körper mit endlich vielen Elementen. Tritt nur in
    Größen \(p^n\) auf (\(p\) prim). Notation \(GF(p^n)\) für
    \emph{Galois-Feld}. Für Reed-Solomon im Datamatrix wird
    \(GF(2^8)\) verwendet}}

\newglossaryentry{galoisfeld}{
  name=Galois-Feld,
  description={Synonym für endlicher Körper, benannt nach Évariste
    Galois (1811--1832). Geschrieben \(GF(p^n)\)}}

\newglossaryentry{charakteristik}{
  name=Charakteristik,
  description={Kleinste positive Zahl \(p\), für die \(p \cdot 1 = 0\)
    gilt. In jedem \(GF(2^n)\) ist die Charakteristik 2 — daraus
    folgt der „Freshman's Dream" \((a+b)^2 = a^2 + b^2\)}}

\newglossaryentry{irreduziblespolynom}{
  name=Irreduzibles Polynom,
  description={Polynom, das nicht in zwei Polynome niedrigeren Grades
    zerlegt werden kann. Analoges Konzept zur Primzahl. Wird gebraucht,
    um \(GF(2^n)\) als Körper zu bauen}}

\newglossaryentry{inverses}{
  name=Multiplikatives Inverses,
  description={Zu einem Element \(a \ne 0\) eines Körpers das
    Element \(a^{-1}\) mit \(a \cdot a^{-1} = 1\). Existiert in
    Körpern für jedes Element außer Null}}

\newglossaryentry{stuetzstelle}{
  name=Stützstelle,
  description={Punkt \(x_i\), an dem ein Polynom \(f\) ausgewertet
    wird, um den Wert \(f(x_i)\) zu erhalten. Bei Reed-Solomon werden
    \(n\) Stützstellen verwendet, von denen \(k\) zur Rekonstruktion
    reichen}}

\newglossaryentry{reedsolomon}{
  name=Reed-Solomon-Code,
  description={Fehlerkorrigierender Code über \(GF(2^n)\). Daten
    werden als Polynom interpretiert und an \(n\) Stützstellen
    ausgewertet; die \(n\) Werte sind das Codewort. Korrigiert bis
    zu \(\lfloor (n-k)/2 \rfloor\) Symbolfehler}}

\newglossaryentry{aushloeschung}{
  name=Auslöschung,
  description={Fehler mit bekannter Position (engl.\ \emph{erasure}).
    Im Gegensatz zum allgemeinen Fehler („wir wissen nicht, welches
    Symbol verfälscht ist") doppelt so günstig zu korrigieren}}

\newglossaryentry{horner}{
  name=Horner-Schema,
  description={Effizientes Verfahren zur Polynom-Auswertung:
    \(a_0 + x(a_1 + x(a_2 + \dots))\). Spart Multiplikationen
    gegenüber der naiven Auswertung}}
